\documentclass{beamer}
\usetheme{Warsaw}

\usepackage[utf8]{inputenc}
\usepackage{fancybox}
\usepackage{multimedia}
\usepackage{subfig}
\usepackage{amsmath}
\usepackage{amssymb}
\usepackage{hyperref}
\hypersetup{
    colorlinks=true,
    urlcolor=blue
}
\usepackage[all]{xy}

\setbeamertemplate{footline}[frame number]

\begin{document}

\title[Stochastik]
{Wahrscheinlichkeitstheorie und Statistik (Stochastik)
\\
\includegraphics[scale=0.5]{img/craps}
}
\subtitle{}
\author[Dr. Johannes Riesterer]
{Dr. rer. nat. Johannes Riesterer}

\date[KPT 2004]{}
\subject{Stochastik}

\frame{\titlepage}

\begin{frame}{\(\sigma\)-Algebra}
\begin{block}{Definition}
Eine Menge \(\mathcal A \subseteq \mathcal P(\Omega)\) hei\ss{}t \(\sigma\)-Algebra auf \(\Omega\), wenn gilt:
\begin{enumerate}
    \item \(\Omega \in \mathcal A\),
    \item aus \(A \in \mathcal A\) folgt \(A^{\mathsf c} \in \mathcal A\),
    \item aus \(A_1,A_2,\dots \in \mathcal A\) folgt
    \[
    \bigcup_{n\in\mathbb N} A_n \in \mathcal A.
    \]
\end{enumerate}
\end{block}
\end{frame}

\begin{frame}{Erzeugte \(\sigma\)-Algebra}
\begin{block}{Definition}
Für ein Mengensystem \(\mathcal G \subseteq \mathcal P(\Omega)\) bezeichnet
\[
\sigma(\mathcal G)
\]
die kleinste \(\sigma\)-Algebra auf \(\Omega\), die \(\mathcal G\) enthält.
\end{block}

\begin{block}{Idee}
Aus den Erzeugern in \(\mathcal G\) werden durch Komplementbildung und abzählbare Vereinigungen alle messbaren Mengen erzeugt.
\end{block}
\end{frame}

\begin{frame}{Existenz und Eindeutigkeit der erzeugten \(\sigma\)-Algebra}
\begin{block}{Satz}
Sei \(\mathcal G\subseteq \mathcal P(\Omega)\) ein Mengensystem.
Dann gibt es genau eine kleinste \(\sigma\)-Algebra auf \(\Omega\), die \(\mathcal G\) enthält.
\end{block}

\begin{block}{Beweisidee}
Betrachte
\[
\mathfrak F
=
\{\mathcal A\subseteq\mathcal P(\Omega):
\mathcal A \text{ ist \(\sigma\)-Algebra und }
\mathcal G\subseteq\mathcal A\}.
\]
Diese Familie ist nicht leer, denn \(\mathcal P(\Omega)\in\mathfrak F\).
\end{block}
\end{frame}

\begin{frame}{Existenz und Eindeutigkeit (Fortsetzung)}
\begin{block}{Beweisidee (Forts.)}
Setze
\[
\sigma(\mathcal G):=\bigcap_{\mathcal A\in\mathfrak F}\mathcal A.
\]
Ein Durchschnitt von \(\sigma\)-Algebren ist wieder eine \(\sigma\)-Algebra und enthält \(\mathcal G\).
Außerdem gilt für jede \(\sigma\)-Algebra \(\mathcal B\) mit \(\mathcal G\subseteq\mathcal B\):
\[
\sigma(\mathcal G)\subseteq\mathcal B.
\]
Also ist \(\sigma(\mathcal G)\) die kleinste solche \(\sigma\)-Algebra.
Die Eindeutigkeit folgt aus der Minimalität.
\end{block}
\end{frame}

\begin{frame}{Borelsche \(\sigma\)-Algebra}
\begin{block}{Definition}
Die borelsche \(\sigma\)-Algebra auf \(\mathbb R\) ist
\[
\mathcal B(\mathbb R)=\sigma(\mathcal O),
\]
wobei \(\mathcal O\) die offenen Mengen in \(\mathbb R\) bezeichnet.
\end{block}

\begin{block}{Wichtige Erzeuger}
Äquivalent gilt
\[
\mathcal B(\mathbb R)=\sigma\bigl(\{(-\infty,a] : a\in\mathbb R\}\bigr).
\]
\end{block}
\end{frame}

\begin{frame}{Messbarkeit reellwertiger Funktionen}
\begin{block}{Definition}
Eine Abbildung
\[
X\colon \Omega \to \mathbb R
\]
hei\ss{}t messbar, wenn für alle \(B\in\mathcal B(\mathbb R)\) gilt:
\[
X^{-1}(B)\in\mathcal A.
\]
\end{block}

\begin{block}{Praktisches Kriterium}
Oft genügt es zu zeigen:
\[
\{\omega\in\Omega : X(\omega)\le a\}\in\mathcal A
\qquad
\text{für alle } a\in\mathbb R.
\]
\end{block}
\end{frame}

\begin{frame}{Maß}
\begin{block}{Definition}
Ein Maß auf einem Messraum \((\Omega,\mathcal A)\) ist eine Abbildung
\[
\mu\colon \mathcal A \to [0,\infty]
\]
mit:
\begin{enumerate}
    \item \(\mu(A)\ge 0\) für alle \(A\in\mathcal A\),
    \item \(\mu(\varnothing)=0\),
    \item für paarweise disjunkte \(A_1,A_2,\dots \in \mathcal A\) gilt
    \[
    \mu\!\left(\bigcup_{n\in\mathbb N} A_n\right)
    =
    \sum_{n\in\mathbb N}\mu(A_n).
    \]
\end{enumerate}
\end{block}
\end{frame}

\begin{frame}{Wahrscheinlichkeitsmaß}
\begin{block}{Definition}
Ein Wahrscheinlichkeitsmaß ist ein Maß \(\mathbb P\) auf \((\Omega,\mathcal A)\) mit der zusätzlichen Eigenschaft
\[
\mathbb P(\Omega)=1.
\]
\end{block}

\begin{block}{Interpretation}
Ein Maß misst allgemein die Größe messbarer Mengen.
Ein Wahrscheinlichkeitsmaß ist der Spezialfall, in dem die Gesamtmasse des ganzen Raums gleich \(1\) ist.
\end{block}
\end{frame}

\begin{frame}{Beispiel: diskretes Wahrscheinlichkeitsmaß}
\begin{block}{Definition}
Für eine endliche Menge \(\Omega=\{\omega_1,\ldots,\omega_n\}\) und Wahrscheinlichkeiten
\(p_1,\ldots,p_n>0\) mit \(\sum_{i=1}^n p_i=1\) definiert man
\[
\mathbb P(A)=\sum_{\omega\in A} p_\omega.
\]
\end{block}

\begin{block}{Spezialfall: Gleichverteilung}
Falls \(p_i=\frac1n\) für alle \(i\), dann gilt
\[
\mathbb P(A)=\frac{|A|}{n}.
\]
\end{block}
\end{frame}

\begin{frame}{Beispiel: Gleichverteilung auf einem Intervall}
\begin{block}{Definition}
Für \([a,b]\subseteq\mathbb R\) mit \(a<b\) ist die Gleichverteilung auf \([a,b]\)
das Wahrscheinlichkeitsmaß mit
\[
\mathbb P([c,d])=\frac{d-c}{b-a}
\qquad
\text{für } a\le c\le d\le b.
\]
\end{block}

\begin{block}{Intuition}
Die Wahrscheinlichkeit ist proportional zur Länge.
Die Dichte ist konstant gleich
\[
\frac1{b-a}.
\]
\end{block}
\end{frame}

\begin{frame}{Zufallsvariable}
\begin{block}{Definition}
Seien \((\Omega,\mathcal A,\mathbb P)\) ein Wahrscheinlichkeitsraum und \((E,\mathcal E)\) ein Messraum.
Eine Abbildung
\[
X\colon \Omega \to E
\]
hei\ss{}t Zufallsvariable, wenn
\[
X^{-1}(B)\in\mathcal A
\qquad
\text{für alle } B\in\mathcal E.
\]
\end{block}
\end{frame}

\begin{frame}{Induzierte Verteilung}
\begin{block}{Definition}
Ist
\[
X\colon (\Omega,\mathcal A,\mathbb P)\to(E,\mathcal E)
\]
eine Zufallsvariable, so ist die induzierte Verteilung von \(X\) das Maß \(\mathbb P_X\) auf \((E,\mathcal E)\) mit
\[
\mathbb P_X(B)=\mathbb P(X\in B)=\mathbb P\bigl(X^{-1}(B)\bigr),
\qquad B\in\mathcal E.
\]
\end{block}
\end{frame}

\begin{frame}{Integration: Grundidee}
\begin{block}{Ziel}
Zu einer messbaren Funktion
\[
f\colon \Omega\to\mathbb R
\]
und einem Maß \(\mu\) möchte man eine Zahl
\[
\int_\Omega f\,d\mu
\]
definieren.
\end{block}

\begin{block}{Lebesgue-Idee}
Nicht das Definitionsgebiet wird zerlegt,
sondern die Funktion wird durch einfache Stufenfunktionen approximiert.
\end{block}
\end{frame}

\begin{frame}{Nichtnegative einfache Funktionen}
\begin{block}{Definition}
Eine nichtnegative einfache Funktion hat die Form
\[
\varphi=\sum_{i=1}^n a_i\,\mathbf 1_{A_i},
\qquad
a_i\ge 0,\ A_i\in\mathcal A.
\]
\end{block}

\begin{block}{Integral}
Für eine solche Funktion setzt man
\[
\int_\Omega \varphi\,d\mu
=
\sum_{i=1}^n a_i\,\mu(A_i).
\]
\end{block}

\begin{block}{Bezug zu Lean}
In mathlib entspricht das dem Begriff der \emph{simple function}.
\end{block}
\end{frame}

\begin{frame}{Warum braucht man Messbarkeit?}
\begin{block}{Idee}
Zur Integration betrachtet man Mengen, auf denen \(f\) bestimmte Werte oder Wertebereiche annimmt, etwa
\[
\{\omega : f(\omega)\ge c\}
\qquad\text{oder}\qquad
\{\omega : a\le f(\omega)<b\}.
\]
\end{block}

\begin{block}{Problem ohne Messbarkeit}
Sind diese Mengen nicht messbar, dann ist \(\mu\) auf ihnen nicht definiert.
Dann kann man weder Stufenfunktionen sinnvoll aufbauen noch ihre Größen messen.
\end{block}

\begin{block}{Merksatz}
Messbarkeit ist genau die Voraussetzung dafür, dass die Niveaumengen von \(f\)
vom Maß erfasst werden.
\end{block}
\end{frame}

\begin{frame}{Wohldefiniertheit für einfache Funktionen}
\begin{block}{Problem}
Eine einfache Funktion kann verschiedene Darstellungen haben:
\[
\varphi=\sum_{i=1}^n a_i\,\mathbf 1_{A_i}
\qquad\text{und}\qquad
\varphi=\sum_{j=1}^m b_j\,\mathbf 1_{B_j}.
\]
\end{block}

\begin{block}{Behauptung}
Der Wert
\[
\int_\Omega \varphi\,d\mu
=
\sum_{i=1}^n a_i\,\mu(A_i)
\]
hängt nicht von der Darstellung ab.
\end{block}
\end{frame}

\begin{frame}{Wohldefiniertheit: Beweisidee}
\begin{block}{Idee}
Man zerlegt den Raum in die Schnitte
\[
C_{ij}=A_i\cap B_j.
\]
Auf jedem \(C_{ij}\) haben beide Darstellungen denselben konstanten Wert,
also liefern sie dieselbe Summe.
\end{block}
\end{frame}

\begin{frame}{Nichtnegative messbare Funktionen}
\begin{block}{Definition}
Sei
\[
f\colon \Omega\to[0,\infty]
\]
messbar.
Dann definiert man
\[
\int_\Omega f\,d\mu
:=
\sup\left\{
\int_\Omega \varphi\,d\mu :
0\le \varphi\le f,\ \varphi \text{ einfach}
\right\}.
\]
\end{block}

\begin{block}{Warum ist das eindeutig?}
Die Menge aller solchen Approximationen ist durch \(f\) und \(\mu\) eindeutig bestimmt.
Also ist auch ihr Supremum eindeutig bestimmt.
\end{block}
\end{frame}

\begin{frame}{Reellwertige Funktionen}
\begin{block}{Zerlegung}
Für eine reellwertige Funktion \(f\) setzt man
\[
f^+(\omega)=\max\{f(\omega),0\},
\qquad
f^-(\omega)=\max\{-f(\omega),0\}.
\]
Dann gilt
\[
f=f^+-f^-,
\qquad
|f|=f^++f^-.
\]
\end{block}

\begin{block}{Definition}
Man setzt
\[
\int_\Omega f\,d\mu
=
\int_\Omega f^+\,d\mu-\int_\Omega f^-\,d\mu,
\]
sofern nicht beide Terme gleichzeitig unendlich sind.
\end{block}
\end{frame}

\begin{frame}{Integrierbarkeit}
\begin{block}{Definition}
Eine reellwertige messbare Funktion \(f\) hei\ss{}t integrierbar, wenn
\[
\int_\Omega |f|\,d\mu<\infty.
\]
\end{block}

\begin{block}{Folge}
Dann sind \(\int f^+\,d\mu\) und \(\int f^-\,d\mu\) endlich,
und damit ist \(\int f\,d\mu\) wohldefiniert und eindeutig.
\end{block}

\begin{block}{Bezug zu Lean}
In mathlib schreibt man daf\"ur \texttt{Integrable f $\mu$}.
\end{block}
\end{frame}

\begin{frame}{Notation in Lean/mathlib}
\begin{block}{Nichtnegatives Integral}
Für Funktionen
\[
f\colon \Omega\to[0,\infty]
\]
verwendet mathlib die Notation
\[
\int^- \omega, f(\omega)\,\partial\mu.
\]
\end{block}

\begin{block}{Gewöhnliches Integral}
Für reellwertige Funktionen schreibt man
\[
\int \omega, f(\omega)\,\partial\mu.
\]
\end{block}
\end{frame}

\begin{frame}{Diskreter Spezialfall}
\begin{block}{Endlicher Fall}
Ist \(\Omega\) endlich und \(\mu\) ein diskretes Maß, dann gilt
\[
\int_\Omega f\,d\mu
=
\sum_{\omega\in\Omega} f(\omega)\,\mu(\{\omega\}).
\]
\end{block}

\begin{block}{Für Wahrscheinlichkeitsmaße}
Ist \(\mu=\mathbb P\), dann ist
\[
\mathbb E[f]=\int_\Omega f\,d\mathbb P
\]
der Erwartungswert.
\end{block}
\end{frame}

\begin{frame}{Messbare Rechtecke}
\begin{block}{Definition}
Für Teilmengen \(A\subseteq E\) und \(B\subseteq F\) hei\ss{}t
\[
A\times B:=\{(x,y)\in E\times F : x\in A,\ y\in B\}
\]
ein Rechteck im Produktraum \(E\times F\).
\end{block}

\begin{block}{Beispiel}
Falls \(E=F=\mathbb R\) und \(A,B\) Intervalle sind, dann ist \(A\times B\subseteq\mathbb R^2\)
ein gewöhnliches Rechteck.
\end{block}
\end{frame}

\begin{frame}{Produkt-\(\sigma\)-Algebra}
\begin{block}{Definition}
Für Messräume \((E,\mathcal E)\) und \((F,\mathcal F)\) ist die Produkt-\(\sigma\)-Algebra
\[
\mathcal E\otimes\mathcal F
=
\sigma\bigl(\{A\times B : A\in\mathcal E,\; B\in\mathcal F\}\bigr).
\]
\end{block}

\begin{block}{Idee}
Sie ist die natürliche \(\sigma\)-Algebra auf dem Produktraum \(E\times F\).
\end{block}
\end{frame}

\begin{frame}{Gemeinsame Verteilung}
\begin{block}{Definition}
Sind
\[
X\colon \Omega\to E,
\qquad
Y\colon \Omega\to F
\]
Zufallsvariablen, so ist
\[
(X,Y)\colon \Omega\to E\times F,
\qquad
\omega\mapsto (X(\omega),Y(\omega)),
\]
eine Zufallsvariable auf \((E\times F,\mathcal E\otimes\mathcal F)\).
\end{block}

\begin{block}{Gemeinsame Verteilung}
Die induzierte Verteilung von \((X,Y)\) hei\ss{}t gemeinsame Verteilung:
\[
\mathbb P_{(X,Y)}(C)=\mathbb P\bigl((X,Y)\in C\bigr),
\qquad
C\in\mathcal E\otimes\mathcal F.
\]
\end{block}
\end{frame}

\begin{frame}{Produktmaß}
\begin{block}{Satz}
Seien \(\mu\) ein Maß auf \((E,\mathcal E)\) und \(\nu\) ein Maß auf \((F,\mathcal F)\).
Dann gibt es genau ein Maß
\[
\mu\otimes\nu
\]
auf \((E\times F,\mathcal E\otimes\mathcal F)\) mit
\[
(\mu\otimes\nu)(A\times B)=\mu(A)\nu(B)
\]
für alle \(A\in\mathcal E\), \(B\in\mathcal F\).
\end{block}
\end{frame}

\begin{frame}{Produktmaß über Schnitte}
\begin{block}{Schnittfunktion}
Für eine messbare Menge \(S\subseteq E\times F\) und \(x\in E\) setze
\[
S_x:=\{y\in F:(x,y)\in S\}.
\]
\end{block}

\begin{block}{Idee}
Dann kann man das Produktmaß durch
\[
(\mu\otimes\nu)(S)
=
\int_E \nu(S_x)\,d\mu(x)
\]
beschreiben.
\end{block}

\begin{block}{Interpretation}
Man misst für jedes \(x\) den vertikalen Schnitt und integriert diese Größen über \(x\).
\end{block}
\end{frame}

\begin{frame}{Produktmaß auf Rechtecken}
\begin{block}{Spezialfall}
Für \(S=A\times B\) gilt
\[
S_x=
\begin{cases}
B,& x\in A,\\
\varnothing,& x\notin A.
\end{cases}
\]
Also
\[
\nu(S_x)=
\begin{cases}
\nu(B),& x\in A,\\
0,& x\notin A.
\end{cases}
\]
\end{block}

\begin{block}{Folgerung}
Daher
\[
(\mu\otimes\nu)(A\times B)
=
\int_E \nu(S_x)\,d\mu(x)
=
\nu(B)\mu(A).
\]
\end{block}
\end{frame}

\begin{frame}{Satz von Fubini}
\begin{block}{Satz}
Seien \((E,\mathcal E,\mu)\) und \((F,\mathcal F,\nu)\) Maßräume und
\[
f\colon E\times F\to\mathbb R
\]
messbar und integrierbar. Dann gilt:
\[
\int_{E\times F} f(x,y)\,d(\mu\otimes\nu)(x,y)
=
\int_E\left(\int_F f(x,y)\,d\nu(y)\right)d\mu(x).
\]
Ebenso gilt
\[
\int_{E\times F} f(x,y)\,d(\mu\otimes\nu)(x,y)
=
\int_F\left(\int_E f(x,y)\,d\mu(x)\right)d\nu(y).
\]
\end{block}
\end{frame}

\begin{frame}{Beweisidee zu Fubini}
\begin{block}{Schritt 1: Indikatorfunktionen}
Für \(f=\mathbf 1_S\) mit messbarem \(S\subseteq E\times F\) gilt
\[
\int_{E\times F}\mathbf 1_S\,d(\mu\otimes\nu)
=
(\mu\otimes\nu)(S),
\]
und nach der Definition des Produktmaßes
\[
(\mu\otimes\nu)(S)=\int_E \nu(S_x)\,d\mu(x).
\]
Da \(\nu(S_x)=\int_F \mathbf 1_S(x,y)\,d\nu(y)\), folgt die Behauptung.
\end{block}

\begin{block}{Schritt 2: Einfache Funktionen}
Mit Linearität folgt der Satz für endliche Summen von Indikatorfunktionen,
also für einfache Funktionen.
\end{block}
\end{frame}

\begin{frame}{Beweisidee zu Fubini: Grenzübergang}
\begin{block}{Schritt 3: Nichtnegative messbare Funktionen}
Man approximiert \(f\ge 0\) von unten durch einfache Funktionen.
Dann erhält man die Aussage mit dem Satz von der monotonen Konvergenz.
\end{block}

\begin{block}{Schritt 4: Allgemeine integrierbare Funktionen}
Für integrierbare \(f\) schreibt man
\[
f=f^+-f^-.
\]
Die Aussage für \(f^+\) und \(f^-\) ist bereits bekannt.
Durch Subtraktion folgt der Satz für \(f\).
\end{block}

\begin{block}{Bemerkung}
Für nichtnegative Funktionen spricht man oft vom Satz von Tonelli,
für integrierbare Funktionen vom Satz von Fubini.
\end{block}
\end{frame}

\begin{frame}{Stochastische Unabhängigkeit}
\begin{block}{Definition}
Zwei Zufallsvariablen
\[
X\colon \Omega\to E,
\qquad
Y\colon \Omega\to F
\]
hei\ss{}en stochastisch unabhängig, wenn für alle \(A\in\mathcal E\), \(B\in\mathcal F\) gilt:
\[
\mathbb P(X\in A,\;Y\in B)=\mathbb P(X\in A)\,\mathbb P(Y\in B).
\]
\end{block}
\end{frame}

\begin{frame}{Charakterisierung der Unabhängigkeit}
\begin{block}{Satz}
Für Zufallsvariablen
\[
X\colon \Omega\to E,\qquad Y\colon \Omega\to F
\]
sind äquivalent:
\begin{enumerate}
    \item \(X\) und \(Y\) sind stochastisch unabhängig,
    \item
    \[
    \mathbb P_{(X,Y)}=\mathbb P_X\otimes\mathbb P_Y.
    \]
\end{enumerate}
\end{block}
\end{frame}

\begin{frame}{Warum ist das plausibel?}
\begin{block}{Beobachtung}
Für alle \(A\in\mathcal E\), \(B\in\mathcal F\) gilt
\[
(X,Y)\in A\times B
\iff
X\in A \text{ und } Y\in B.
\]
\end{block}

\begin{block}{Daher}
\[
\mathbb P_{(X,Y)}(A\times B)
=
\mathbb P(X\in A,\;Y\in B).
\]
Also kodiert die gemeinsame Verteilung auf Rechtecken genau die gemeinsamen Auftretenswahrscheinlichkeiten.
\end{block}
\end{frame}

\begin{frame}{Beweisidee: \(1\Rightarrow 2\)}
\begin{block}{\(1\Rightarrow 2\)}
Aus der Unabhängigkeit folgt auf jedem Rechteck
\[
\mathbb P_{(X,Y)}(A\times B)
=
\mathbb P_X(A)\mathbb P_Y(B)
=
(\mathbb P_X\otimes\mathbb P_Y)(A\times B).
\]
Da Rechtecke die Produkt-\(\sigma\)-Algebra erzeugen, folgt
\[
\mathbb P_{(X,Y)}=\mathbb P_X\otimes\mathbb P_Y.
\]
\end{block}
\end{frame}

\begin{frame}{Beweisidee: \(2\Rightarrow 1\)}
\begin{block}{\(2\Rightarrow 1\)}
Gilt umgekehrt
\[
\mathbb P_{(X,Y)}=\mathbb P_X\otimes\mathbb P_Y,
\]
so erhält man auf Rechtecken sofort
\[
\mathbb P(X\in A,\;Y\in B)=\mathbb P(X\in A)\mathbb P(Y\in B).
\]
Also sind \(X\) und \(Y\) unabhängig.
\end{block}
\end{frame}

\end{document}